\documentclass[amssymb,amsmath,12pt,letterpaper]{revtex4}

\renewcommand{\rmdefault}{phv} % Arial
\renewcommand{\sfdefault}{phv} % Arial

\usepackage{graphicx}
\usepackage[spanish]{babel}
\usepackage[utf8]{inputenc}
\usepackage{tabulary}

\linespread{1.3} % one and a half spacing

\begin{document}

% \begin{minipage}[center]{1.2\textwidth}
  
%   \vspace{1cm}
%   \begin{tabular}{m{0.1\textwidth}m{0.7\textwidth}}
%     \includegraphics[width=0.1\textwidth]{logo_UdeA} & \centerline{\textit{Instituto de Física, Universidad de Antioquia.}}
%   \end{tabular}

% \end{minipage}


\title{\textsc{Solución de la ecuación de Laplace a través del metodo de relajación.}}

\author{Sebastián Duque Mesa.\footnote{Correo Electrónico: sduquemesa@gmail.com}}       %autores

\affiliation{Instituto de Física, Universidad de Antioquia, A.A. 1226, Medellín, Colombia.}


\begin{abstract}
                                
  \vspace{0.5cm}
  \begin{picture}(1,1)(1,1) \put(1,10){\line(1,0){380}} \end{picture}                   
  \begin{center}\textbf{\textsc{Resumen.}}\end{center}

  %---------------------------------------------------------
  Es posible encontrar una solución aproximada a la ecuación de Laplace utilizando únicamente aritmética.

  Si se realizan expansiones de Taylor para una función de potencial $\phi(x,y,z)$ alrededor de un punto $(x_0,y_0,z_0)$ es posible demostrar que el potencial en este punto es el promedio de los potenciales en los puntos localizado sobre una esfera de radio $\delta$ centrada en este.

  Para ilustrar, utilizando el resultado anterior, calculamos el potencial dentro de un capacitor coaxial infinito de sección transversal cuadrada y, a partir de los resultados, se dibujan lineas equipotenciales dentro del capacitor.
  %---------------------------------------------------------
                                
  \vspace{0.8cm}
  \begin{picture}(1,1)(1,1) \put(1,10){\line(1,0){380}} \end{picture}
  
\end{abstract}

\vspace{0.5cm}
\maketitle

Para la realización ponencia se solicita la disposición de un Video Beam y tablero para marcador.

Linea de trabajo de la ponencia: Física pura.

\end{document}
